\section{An exact hidden-model memory frontier}\label{sec:r79-exact-hidden}

The model family in this example consists of two scalar instruments.
The model $\theta\in\{0,1\}$ stays fixed during the run and is not observed.
There are two commands, $a$ and $b$, one failure report, unit call cost,
and compulsory stopping at success.  In model $0$, command $a$ succeeds
with probability one and $b$ fails with probability one; in model $1$
the roles are reversed.  The controller knows its sampled command,
and may use that command together with the failure report in its update.
It is autonomous, has a prescribed pure initial live label, and has no
uncounted clock or persistent random seed.

\begin{proposition}[Golden-ratio memory cost]
\label{prop:r79-golden-ratio-memory}
Writing $\varphi=(1+\sqrt5)/2$, the worst-model expected total cost obeys
\[
 V_1=2,\qquad V_2=\varphi,\qquad V_M=\frac32\quad(M\ge3).
\]
The last value is also the optimum among unrestricted history-dependent
controllers.  For the two-label optimum, one label is unreachable and
improper in one model, although the prescribed-start controller is proper
in both models.  Thus model-dependent reachable supports cannot in
general be replaced by a full-space resolvent or full-space Lyapunov
inequality.
\end{proposition}
\begin{proof}
At a live label let $r$ denote the probability of command $a$.  A single
label has expected costs $1/r$ and $1/(1-r)$, interpreted as infinite at
the corresponding endpoints.  Their maximum is minimized at $r=1/2$
and equals $2$.

For two labels let the command probabilities be $r,s$, with initial row
$r$.  At every later call the success probability in model $0$ is at most
$\max(r,s)$ and in model $1$ at most $1-\min(r,s)$.  A process whose
conditional success probability is at most $h>0$ has expected number of
remaining calls at least $1/h$: its survival probabilities are at least
$(1-h)^n$, by induction.  If $h=0$, that expectation is infinite.
Conditioning on the first call therefore gives
\[
 C_0\ge1+\frac{1-r}{\max(r,s)},\qquad
 C_1\ge1+\frac{r}{1-\min(r,s)},
\]
whenever the displayed terms have positive numerator; endpoint cases
with zero numerator have the evident zero continuation contribution.
If $r\ge s$, these bounds imply
\[
 \max(C_0,C_1)\ge\max\{1/r,1+r\}\ge\varphi.
\]
If $r\le s$, they imply
\[
 \max(C_0,C_1)\ge\max\{2-r,1/(1-r)\}\ge\varphi.
\]
The boundary cases satisfy the same lower bound by direct inspection.
Both one-variable minima follow by equating their decreasing and
increasing entries.

For attainment take $r=1/\varphi=\varphi-1$ and $s=0$.  From label $1$,
a failed $a$ command writes label $2$, whose command is always $b$,
while a failed $b$ command retains label $1$.  At label $2$, every
failure retains label $2$.  In model $0$, only $b$ can fail, so label
$2$ is never reached and the expected cost is $1/r=\varphi$.  In model
$1$, the first call either succeeds on $b$ or fails on $a$ and is followed
by one certain success, giving cost $1+r=\varphi$.  This computes the
actual cost without assuming properness of the unused model-$0$ label.

For an arbitrary history-dependent controller, let $r$ be its first-call
probability of $a$.  A failed first call needs at least one additional
unit-cost call before success.  The two model costs are therefore at least
$2-r$ and $1+r$, whose maximum is at least $3/2$.  Three counted labels
attain the bound: use an initial $1/2$--$1/2$ row and two continuing rows,
one always $a$ and one always $b$.  A failed $a$ command writes the pure-$b$
label and a failed $b$ command writes the pure-$a$ label.  The next call
succeeds surely.  Duplicating labels gives the assertion for every $M\ge3$.

Finally, under the two-label optimizer the scalar failure operator in
model $0$ leaves label $2$ invariant, since its $b$ command always fails.
A full-space cost supersolution would require $W_{0,2}\ge1+W_{0,2}$,
which is impossible.  The valid reachable projections are
$P_0=(1,0)$ and $P_1=(1,1)$; the corresponding cost certificates are
$W_0=(\varphi,0)$ and $W_1=(\varphi,1)$.  They obey the support-restricted
Bellman equalities and the same initial bound $W_{\theta,1}=\varphi$.
\end{proof}

The pure initial label matters.  If an uncounted randomized initial
choice of labels were free, two pure command rows would already attain
$3/2$ by an equal initial mixture and switching after the first failure.
That is a different memory convention.  Likewise a continuously sampled
persistent seed or externally supplied time index would change the
controller class.


